\section{Integration with Prior Results}
\label{sec:integration}

This section records how controlled relaxation extends the
existing GFM machinery: it preserves self-balancing, resolves
inter-agent risk disagreement, and extends the multi-channel
detection architecture with a sixth, actively generated channel.

\subsection{Preservation of self-balancing}

\begin{proposition}[Controlled Relaxation Preserves Self-Balancing]
\label{prop:preserve_balance}
The self-balancing property of $\volL$
(\citep[Theorem~1]{lasser2026horizon}, self-balancing of the
discounted value under capability dynamics) is preserved under
controlled relaxation:
\begin{itemize}
\item[(a)] During the test window ($\Xirestrict$ lifted within
  $\scope$), the capability tuple $(d_1, \ldots, d_k)$ is
  exercisable, which by weak-monotonicity axiom~M4
  of~\citep[Proposition~1]{lasser2026poset} means $\volL$ is at
  least as large as under restriction, minus the possible
  damage bounded by $\damagebound$.  The net effect on $\volL$
  is bounded by
  $[-\damagebound, +\Delta_{\mathrm{unrestrict}}]$ where
  $\Delta_{\mathrm{unrestrict}} \geq 0$ is the $\volL$ gain from
  lifting the restriction.
\item[(b)] After the test, the updated risk estimate
  $\pR(t + \testdur)$ more accurately reflects the true risk, so
  the restriction
  criterion~\citep[Proposition~1]{lasser2026horizon} is applied
  against a more accurate estimate.  If the risk was genuinely
  low, the restriction is lifted and $\volL$ permanently
  increases; if the risk was genuinely high, the restriction is
  maintained with higher confidence and $\volL$ is unchanged
  (the restriction was already in place).
\item[(c)] Over repeated tests
  (Theorem~\ref{thm:convergence}), the risk estimates converge
  to truth, and the restriction landscape converges to the
  optimal restriction set---the set of tuples that are
  genuinely risky.  This is a stronger self-balancing guarantee
  than the static version: not only does $\volL$ self-balance
  against perturbations, but the restriction landscape itself
  converges to the truth-tracking configuration.
\end{itemize}
\end{proposition}

\begin{proof}[Proof sketch]
Parts~(a) and (b) are direct consequences of
Theorem~\ref{thm:damage_bound} and the criterion
of~\citep[Proposition~1]{lasser2026horizon} applied to the
updated $\pR$.  Part~(c) follows from
Theorem~\ref{thm:convergence}: under
Assumptions~\ref{ass:T1}--\ref{ass:T5}, the posterior converges
almost surely to the correct hypothesis; iterating this over
all restricted tuples yields convergence of the restriction
landscape.  Technical details in
Appendix~\ref{app:proof_preserve_balance}.
\end{proof}

\subsection{Empirical arbitration across disagreeing agents}

\begin{proposition}[Risk-Claim Consensus via Empirical Arbitration]
\label{prop:empirical_arbitration}
When multiple agents disagree on the risk of a capability tuple
($a_j$ assigns $\pR^{(j)} \gg \pR^{(k)}$ for agents $j, k$),
controlled relaxation provides an empirical arbitration
mechanism.
\begin{itemize}
\item[(a)] Under max-aggregation of the risk
  estimates~\citep{lasser2026horizon}, agent $a_j$'s high
  estimate locks the capability.  Under log-odds
  aggregation~\citep[Definition~5]{lasser2026scm}, the
  disagreement produces a moderate aggregate that may or may
  not trigger restriction.
\item[(b)] In either case, if the capability is restricted,
  controlled relaxation generates evidence that updates both
  agents' estimates toward the truth
  (Theorem~\ref{thm:convergence}).  After sufficient tests, the
  agents' posteriors converge, and the disagreement is resolved
  by data rather than by aggregation-rule choice.
\item[(c)] The convergence rate depends on each agent's prior
  weight (how strongly they hold their belief), but
  Assumption~\ref{ass:T4} (prior support) and the KL-positivity
  of Theorem~\ref{thm:convergence}(d) guarantee convergence
  regardless of prior: even a strongly held false risk claim is
  eventually overwhelmed by repeated controlled-relaxation
  evidence.
\end{itemize}
\end{proposition}

\begin{proof}[Proof sketch]
Part~(a) restates the aggregation-rule behaviour
of~\citep{lasser2026horizon, lasser2026scm}.  Parts~(b)
and~(c) follow directly from Theorem~\ref{thm:convergence}
applied independently to each agent's posterior, using
identical evidence streams.
\end{proof}

\subsection{A sixth detection channel}

\begin{lemma}[Controlled Relaxation as a Sixth Detection Channel]
\label{lem:sixth_channel}
The controlled relaxation test adds a sixth channel to the
five-channel attribution
architecture~\citep[Table~1]{lasser2026scm}:
\begin{table}[h]
\centering
\small
\begin{tabular}{@{}p{4.0cm}p{4.2cm}p{4.2cm}@{}}
\toprule
\textbf{Channel} & \textbf{Source} & \textbf{Signal} \\
\midrule
1. Capability gating
  & Exercise outcomes
  & Binary (pass/fail) \\
2. Causal attribution
  & SCM interventions
  & Continuous ($\Delta \volL$) \\
3. Behavioural residual
  & Agent policy
  & Continuous (deviation) \\
4. SPRT detection
  & Sequential test
  & Binary (accept/reject) \\
5. Correlational co-occurrence
  & Pattern matching
  & Continuous (correlation) \\
\textbf{6. Controlled relaxation}
  & \textbf{Direct exercise under monitoring}
  & \textbf{Binary (contraction / no contraction) plus
    continuous ($\volL$ delta)} \\
\bottomrule
\end{tabular}
\caption{Extended detection architecture, adding controlled
  relaxation as a sixth channel.  Rows 1--5
  from~\citep[Table~1]{lasser2026scm}.}
\label{tab:sixth_channel}
\end{table}

\noindent The sixth channel is distinct from the others in that
it is \emph{actively generated} rather than passively observed.
The framework designs the observation (by constructing the
scope and lifting the restriction), rather than waiting for the
observation to arrive.  This active-experiment property makes
it the only channel that can generate evidence \emph{against} a
risk claim under the verification asymmetry.
\end{lemma}

\begin{proof}[Proof sketch]
The sixth channel's signal is
$O_n \in \{0, 1\}$ plus the measured
$\Delta \volL$ during the test window.  Its isolation and
independence properties inherit from
Assumptions~\ref{ass:D1}--\ref{ass:D2} combined with channel
isolation~\citep[Proposition~5]{lasser2026exo}: the active-exercise
signal is cross-substrate-isolated from the
agent-claim channel because it arises from direct poset-state
measurement rather than agent-mediated report.
\end{proof}

\begin{remark}[Test outcome vs.\ test occurrence]
\label{rem:outcome_vs_occurrence}
The independence claim above holds at the level of test
\emph{outcome} given test \emph{occurrence}: $O_n$ is
independent of agent reports conditional on the test having
been run.  The test \emph{occurrence} itself, however, is a
function of the meta-policy (Section~\ref{sec:meta_policy}),
which reads from agent reports to decide whether a restriction
is stale and whether the cost--benefit favours testing.  For
the log-odds aggregation
of~\citep[Definition~5]{lasser2026scm}, the distinction matters
only if the aggregation conditions on test-occurrence when
combining evidence with other channels; the protocol does not.
The joint-factorization assumption
of~\citep{lasser2026exo} is satisfied in the
outcome-conditional-on-occurrence sense, not in the
occurrence-marginal sense.
\end{remark}

\begin{remark}[Trust-weight composition with detection integrity]
\label{rem:trust_weight_composition}
The trust weight of $1$ on the CR channel above presupposes
that the detection apparatus is correctly calibrated and the
scope correctly constructed.  Both are themselves subject to
integrity constraints.  A robust composition caps the
CR-channel trust weight by the integrity guarantee of the
underlying detection apparatus:
$w_{\mathrm{CR}} = \min(1, w_{\mathrm{integrity}})$, where
$w_{\mathrm{integrity}}$ is obtained from the cryptographic
commitment apparatus of~\citep{lasser2026exo} when one applies,
or from the institutional integrity assessment otherwise.  Under
this composition, integrity failures in the detection apparatus
propagate to the risk estimate with bounded weight rather than
full weight; weight-$1$ trust on the CR channel is the special
case $w_{\mathrm{integrity}} \geq 1$, which is rarely defensible
in practice.  The composition makes explicit that the CR
channel's trust inherits from the trust mechanism rather than
being independent of it.
\end{remark}
